\documentclass[12pt]{ctexart}
\usepackage{amsmath}
\title{风光储能协同优化模型}
\author{参赛队}
\begin{document}
\maketitle
\begin{abstract}
本文针对风光储能协同调度问题建立混合整数规划模型，目标为最小化 144 个时段的购电成本，
并在功率平衡、充放电互斥与容量边界约束下求解。
\end{abstract}
\section{问题重述}
题目要求给出逐时段计划购电曲线，并满足供需平衡与 SOC 递推约束。
\section{模型建立}
目标函数为
\begin{equation}
\min \sum_{t=1}^{144} c_t p_t
\label{eq:objective}
\end{equation}
约束包括功率平衡、充放电互斥与容量边界，如式~\eqref{eq:objective} 所示。
\section{求解结果}
求解得到日内购电成本较基准方案下降 3.1 个百分点。
\section{检验与推广}
对光伏出力做正负十个百分点灵敏度分析，结论保持稳定。
\end{document}